\documentclass{beamer}
% \documentclass[handout]{beamer}
\setbeamertemplate{navigation symbols}{}
\usepackage{beamerthemeshadow}
\usepackage[ngerman]{babel}	% Spracheinstellungen
\usepackage[utf8]{inputenc}	% Zeichencodierung
\usepackage[T1]{fontenc}
\beamersetuncovermixins{\opaqueness<1>{25}}{\opaqueness<2->{15}}


% --------------------------- Alles fuer die Code-Listings -----------------------------
\usepackage{xcolor}
\definecolor{ListingBackground}{rgb}{0.9,0.9,0.9}
\definecolor{DarkBlue}{rgb}{0,0.08,0.45}
\definecolor{NiceRed}{rgb}{0.89, 0.0, 0.52}
\definecolor{CodeGreen}{rgb}{0.0, 0.49, 0.08}
\usepackage{listings}		% Code einbinden
\usepackage{green}
% \usepackage{blue}

\setbeamertemplate{footline}
{%
  \leavevmode%
  \hbox{\begin{beamercolorbox}[wd=.5\paperwidth,ht=2.5ex,dp=1.125ex,leftskip=.3cm plus1fill,rightskip=.3cm]{author in head/foot}%
    \usebeamerfont{author in head/foot}\insertshortauthor
  \end{beamercolorbox}%
  \begin{beamercolorbox}[wd=.5\paperwidth,ht=2.5ex,dp=1.125ex,leftskip=.3cm,rightskip=.3cm plus1fil]{title in head/foot}%
    \usebeamerfont{title in head/foot}\insertshortinstitute
%     \includegraphics[width=20px]{univentionklein}
  \end{beamercolorbox}}%
  \vskip0pt%
}

\titlegraphic{\includegraphics[width=0.3\columnwidth]{technik} \hskip0.5cm \includegraphics[width=0.3\columnwidth]{univentionkomplett}} 
% 	Fachbereich Technik \\
% 	Abteilung Elektrotechnik und Informatik
% \setbeamertemplate{enumerate item}{\color{technikdunkel}\insertenumlabel.}
\setbeamertemplate{enumerate item}{\begin{overpic}[scale=.4,unit=1mm]%
{kreis.pdf}
\put(25,20){\ttfamily\fontsize{9pt}{9pt}\selectfont\color{white}\insertenumlabel}
\end{overpic}}
% \ttfamily\fontsize{9pt}{10pt}\selectfont
% \ttfamily\fontsize{9pt}{10pt}\selectfont
%\setbeamertemplate{itemize item}{\color{technikdunkel}\ttfamily\fontsize{90pt}\selectfont.}

\begin{document}
\title{Effiziente Speicherung virtueller Festplatten mit bestehender OpenSource-Software}  
\author{Bastian de Groot}
\institute{Hochschule Emden-Leer}
\date{\today} 
%\definecolor{Name}{System}{Zusammensetzung}


\begin{frame}
\titlepage
\end{frame} 

\section{Virtualisierung}
\begin{frame}
\frametitle{(Betriebsystem-)Virtualisierung}
\fontsize{15pt}{25pt} \selectfont
\begin{itemize}
 \item Verteilung von Hardwareressourcen auf virtuelle Maschinen \pause  %Betriebsystem-Instanzen
 \item Virtuelle Maschinen stellen nachgebildete Rechner dar, die selbst (in der Regel) keinen Unterschied zu echter Hardware feststellen können
\end{itemize}
\end{frame}

\section{Copy-on-Write}

\begin{frame}
\frametitle{Vorraussetzungen effizienter Speicherung}
\fontsize{20pt}{25pt} \selectfont
\begin{itemize}
 \item Viele virtuelle Maschinen \pause
 \item Große Ähnlichkeit untereinander 
%  \item die virtuellen Maschinen sollen schnell verfügbar sein
\end{itemize}

\end{frame}

\begin{frame}
\frametitle{Virtualisierte Hardware} 
\fontsize{20pt}{25pt} \selectfont
\begin{itemize}
 \only<1->{ \item \includegraphics[width=20px]{cpu} \hskip0.2cm Prozessor} 
 \only<2->{ \item \includegraphics[width=20px]{festplatte} \hskip0.2cm Festplatte} 
 \only<3->{ \item \includegraphics[width=20px]{cd} \hskip0.2cm CD-Laufwerk} 
 \only<4->{ \item \includegraphics[width=20px]{monitor} \hskip0.2cm Monitor} 
 \only<5->{ \item \includegraphics[width=20px]{maus} \hskip0.2cm Maus} 
 \only<6->{ \item \includegraphics[width=20px]{tastatur} \hskip0.2cm Tastatur} 
\end{itemize}
\end{frame}

\begin{frame}
\frametitle{Kritische Hardware für die Speicherung} 
\fontsize{20pt}{25pt} \selectfont
\begin{itemize}
 \only<1-1>{ \item \includegraphics[width=20px]{cpu} \hskip0.2cm Prozessor} 
 \only<1-2>{ \item \includegraphics[width=20px]{festplatte} \hskip0.2cm Festplatte} 
 \only<1-2>{ \item \includegraphics[width=20px]{cd} \hskip0.2cm CD-Laufwerk} 
 \only<1-1>{ \item \includegraphics[width=20px]{monitor} \hskip0.2cm Monitor} 
 \only<1-1>{ \item \includegraphics[width=20px]{maus} \hskip0.2cm Maus} 
 \only<1-1>{ \item \includegraphics[width=20px]{tastatur} \hskip0.2cm Tastatur} 
\end{itemize}
%  \only<4>{\includegraphics[width=150px]{festplatte} \\ Festplatte}
\end{frame}

\begin{frame}
\frametitle{Kritische Hardware für die Speicherung} 
\begin{center}
 \includegraphics[width=150px]{festplatte} \\ Festplatte
\end{center}

\end{frame}

\begin{frame}
\frametitle{Kritische Hardware für die Speicherung} 


\begin{columns}
\column{3.8cm} 
\begin{center}
 \includegraphics[width=100px]{festplatte} \\ Festplatte
\end{center}
\column{6cm} 
\begin{block}{Speicherungsformen}
 \begin{itemize}
 \item Festplattenpartition \pause
 \item Rohe Abbildungsdatei (Raw-Image) \pause
 \item Viele Abbildungsformate der Virtualisierungstechnologien: vmdk, qcow2, vdi, vhd, ...
 \end{itemize}
\end{block}
\end{columns}
\end{frame}

\begin{frame}
\frametitle{Herkömmliche/schlechte Lösung} 
\includegraphics<1->[height=100px]{copyonwrite1}
\includegraphics<2->[height=100px]{copyonwrite2} 
\includegraphics<3->[height=100px]{copyonwrite3}
\end{frame}

\begin{frame}
\frametitle{Gute Lösung - Copy-on-Write} 
\includegraphics<1->[height=100px]{copyonwrite_1}
\includegraphics<2->[height=100px]{copyonwrite_2} 
\includegraphics<3->[height=100px]{copyonwrite_3}
\end{frame}

\begin{frame}
\frametitle{Gute Lösung - Copy-on-Write} 
\includegraphics[height=100px]{copyonwrite_1}
\includegraphics[height=100px]{copyonwrite_2} 
\includegraphics[height=100px]{copyonwrite_3}
\begin{center}
\vskip0.3cm
\includegraphics[width=60px]{Smilie}
\end{center}
\end{frame}

\begin{frame}
\frametitle{Fazit Copy-on-Write}

\begin{block}{Vorteile}
 \begin{itemize}
  \item Sehr schnelles Klonen von virtuellen Maschinen
  \item Sehr effiziente Speicherung
 \end{itemize}

\end{block}

\begin{alertblock}{Zu beachten}
 \begin{itemize}
  \item Inkompatibilitäten zwischen Copy-on-Write Imageformaten
 \end{itemize}
\end{alertblock}

\end{frame}


\section{Verteilung}
\begin{frame}
\frametitle{Verteilung der Master-Images - Vorhandene Lösungen} 
\fontsize{20pt}{30pt} \selectfont
\begin{itemize}
 \item BitTorrent
 \item Multicast
 \item NFS
\end{itemize}
\end{frame}

\begin{frame}
\frametitle{Multicast} 
\includegraphics[width=300px]{multicast}
\end{frame}

\begin{frame}
\begin{columns}
\column{3.8cm} 
\frametitle{Multicast} 
\includegraphics[width=120px]{multicast}
\column{6cm} 
\begin{block}{Vorteile}
\begin{itemize}
 \item Sehr hohe Geschwindigkeit durch Parallelität \pause
\end{itemize}
\end{block}

\begin{alertblock}{Nachteile}
\begin{itemize}
 \item Hohe Netzwerklast während der Verteilung \pause
 \item Geschwindigkeitseinbruch bei heterogener Umgebung oder schlechten Netzanbindungen
\end{itemize}
\end{alertblock}
\end{columns}
\end{frame}

\begin{frame}
\frametitle{NFS} 
\includegraphics[width=300px]{nfs}
\end{frame}

\begin{frame}
\begin{columns}
\column{3.8cm} 
\frametitle{NFS} 
\includegraphics[width=120px]{nfs}
\column{6cm} 
\begin{block}{Vorteile}
\begin{itemize}
 \item Geringer Administrationsaufwand \pause
 \item Sofortige Verfügbarkeit \pause
\end{itemize}
\end{block}

\begin{alertblock}{Nachteile}
\begin{itemize}
 \item Permanent hohe Netzwerklast, bei mehreren gestarteten VMs \pause
 \item Keine Lastenverteilung
\end{itemize}
\end{alertblock}
\end{columns}
\end{frame}

\begin{frame}
\frametitle{BitTorrent} 
\includegraphics[width=300px]{bittorrent}
\end{frame}

\begin{frame}
\frametitle{BitTorrent} 
\begin{columns}
\column{3.8cm} 
\includegraphics[width=120px]{bittorrent}
\column{6cm} 
\begin{block}{Vorteile}
\begin{itemize}
 \item Hohe Skalierbarkeit \pause 
%  \item Sehr effizient auch in heterogenen Umgebungen \pause
 \item Hohe Ausfallsicherheit \pause
\end{itemize}
\end{block}

\begin{alertblock}{Nachteile}
\begin{itemize}
  \item Leicht verringerte Geschwindigkeit bei asymmetrischer Upload- und Downloadgeschwindigkeit
\end{itemize}
\end{alertblock}
\end{columns}
\end{frame}

\begin{frame}
\frametitle{Fazit} 
\fontsize{20pt}{25pt} \selectfont
\begin{itemize}
  \item Multicast ungeeignet, da alle Daten ankommen müssen \pause
  \item NFS ungeeignet, da es keine Lastenverteilung hat \pause
  \item BitTorrent ist die beste Lösung
\end{itemize}
\end{frame}


\section{Softwarelösung}

\begin{frame}
\frametitle{Softwarelösung} 
\includegraphics[width=300px]{Kommunikation}

\end{frame}

\begin{frame}
\begin{columns}
\column{3.8cm} 
\frametitle{Softwarelösung} 
\includegraphics[width=120px]{Kommunikation}
\column{6cm} 
\begin{block}{Verwendungszweck}
\begin{itemize}
 \item Klonen virtueller Maschinen \pause
 \item Verteilen von Master-Images im Netzwerk
\end{itemize}
\end{block}
\end{columns}
\end{frame}

\begin{frame}
\frametitle{Vorführung} 
\begin{center}
 \includegraphics[width=150px]{video} 
\end{center}

\end{frame}


\begin{frame}
\frametitle{Ende}
\fontsize{70pt}{70pt} \selectfont Fragen
\end{frame}

\end{document}